% ============ 第一部分：摘要后的第 1--3 章 ============

\section{引论：一道题的四种问法与一条被选中的路}

\subsection{官方提法与四个备选}

克雷数学研究所在 2000 年把 \NS{} 方程的存在性与光滑性问题列为七个千禧年问题之一。查尔斯·费弗曼（Charles Fefferman）撰写的官方提法\cite{clay}并没有给出单一命题，而是并列了四条互不依赖的路径，声明“证明其中任何一条即可”。其中两条是存在性陈述：三维空间中（备选 A）与周期盒 $\mathbb{T}^{3}$ 中（备选 B），从任意光滑、速降且无散度的初值出发，方程都有整体光滑解，且动能一致有界；另外两条是反例陈述：空间情形（备选 C）与周期情形（备选 D）允许求解者同时指定一个光滑初值和一个光滑外力，只要外力及其全部导数满足速降条件，使得不存在具有一致有界动能的整体光滑解。

\begin{keybox}
四条备选之间不是逻辑否定的关系。备选 C 允许在外力容许的条件下取反例，它不是备选 A 的否定：A 的否定应是“某个无外力初值没有整体光滑有界能量解”，而 C 的反例额外挂上了一项光滑外力。因此，构造出 C 只说明“在某些光滑推动下方程会失稳”，并不裁决“没有任何外力时方程是否永远光滑”。
\end{keybox}

官方文本对此的解释只有一句话，即这四条“给求解者合理的余地，同时保留问题的核心”。这句话本身就承认了：外力是问题的一部分，而不是外挂的技巧。

\subsection{2026 年 9 月发生了什么}

按 OpenAI 官方说明\cite{openai-blog}与其手稿\cite{openai-paper}，事件的时间轮廓大致如下：公司自 8 月 28 日起训练一个内部模型；9 月 1 日听闻“两个千禧年问题已被解决”的传闻，遂组织大规模智能体系统对若干公开难题发起搜索；先由约一百个智能体用约五十小时解决无外力三维 Euler 方程的爆破问题，随后把资源集中到 \NS{}；约一万个并发智能体在 88 小时后（9 月 5 日）给出 \NS{} 爆破解，Lean 形式化与验证再由 GPT-6 Astra 花费约十七小时完成。9 月 8 日，随官方说明一同发布的还有 166 页手稿《Finite time blowup for Navier--Stokes》与公开的 Lean 仓库\cite{openai-lean}；同日，纽约大学柯朗研究所的 Tristan Buckmaster 与 Anthropic 的 Levent Alpöge 也公开了他们在外力作用下的 IPM、Boussinesq 与三维 Euler 爆破结果，并对“抢发”提出异议。9 月 11 日，克雷研究所表态称该问题“看上去已被解决”，但强调评估与授信的流程“刻意不匆忙”\cite{clay-news}。

\begin{table}[H]
\centering
\small
\caption{2026 年 9--10 月关键节点（据公开来源整理）}
\begin{tabular}{@{}>{\raggedright\arraybackslash}p{2.6cm}>{\raggedright\arraybackslash}p{11.4cm}@{}}
\toprule
\hdrow \textbf{日期} & \textbf{事件} \\
\midrule
8/15--8/22 & Alpöge--Buckmaster 获得光滑外力下 IPM/Boussinesq/三维 Euler 的爆破结果，并完成 Lean 验证。 \\
9/1 & OpenAI 听闻传闻，启动多智能体战役。 \\
9/5--9/6 & 约一万并发智能体经 88 小时产出 \NS{} 爆破证明；Lean 形式化再耗约 17 小时。 \\
9/7 & Tao 发表博文，梳理 Alpöge--Buckmaster 的方法谱系，指向 Córdoba--Martínez-Zoroa。 \\
9/8 & OpenAI 发布手稿与 Lean 仓库；Buckmaster 同日发布“抢发”声明。 \\
9/11 & 克雷研究所表态“看上去已被解决”，评估“刻意不匆忙”。 \\
9/17 & Constantin--Ignatova--Vicol 上传刚性定理预印本；Gowers 撰文解释为何未签署菲尔兹奖得主公开信。 \\
10/6 & Bastounis--Circelli--Hansen 上传关于自动形式化非语义忠实的预印本，以引理 8.6 为例。 \\
10/10 & 媒体报道 OpenAI 已放出 722 篇 AI 生成数学手稿，学界反弹加剧。 \\
\bottomrule
\end{tabular}
\end{table}

\subsection{本文的范围、立场与方法}

本文是一篇基于公开材料的综述，而不是对该证明的独立复核。凡涉及“证明是否正确”的判断，本文一律转述公开文献并标明其作者的立场与保留意见；凡涉及技术内容的陈述，则以手稿原文、Lean 仓库与已公开的独立文献为准，并给出可核验的出处编号。这样做有两条理由。其一，166 页手稿加六十余万行 Lean 代码，其完全消化需要相当的时间，任何在数日内给出的“已确认/已推翻”结论都不诚实。其二，本文关心的问题不是“是非”，而是“这项工作的数学思想从哪里来、用了什么技术、这些技术和它们的验证处在什么位置”——这类问题可以依据公开材料稳妥地回答。

\begin{principlebox}
\textbf{本文主张的主线}\quad 这项工作的数学内容，可以压缩成一句话：先写下想要的奇异解，再让外力恰好等于残差；全部的困难都落在“让残差连同它的一切导数光滑地穿过奇点”上。围绕这一句话，才长出涡旋的标度、脉冲的通量、逐阶的修正与环面上的乘子估计。理解了这句话，才知道后面每一块技术各自在做什么。
\end{principlebox}

\subsection{记号约定}

全文用 $u(x,t)$ 表示速度场，$p$（或 $\pi$）表示压力，$f$ 表示外力，$\nu>0$ 为黏性系数。三组方程分别写作
\begin{align}
\text{（\NS{}）}\quad & \partial_{t}u+(u\cdot\nabla)u=-\nabla p+\nu\Delta u+f,\qquad \nabla\cdot u=0;\\
\text{（Euler）}\quad & \partial_{t}u+(u\cdot\nabla)u=-\nabla p+f,\qquad \nabla\cdot u=0.
\end{align}
柱坐标取 $x=(r\cos\theta,r\sin\theta,z)$，$u_r,u_\theta,u_z$ 为对应的三个分量。记奇点时间前的剩余时间为 $\tau=1-t$，径向与轴向特征长度分别记为 $\ell_r,\ell_z$。$\|u\|_{L^{\infty}}$ 指速度的 $L^{\infty}$ 范数，$\|u\|_{L^{2}}$ 指点动能的平方根。书写上，$\mathcal{P}u$ 表示速度关于 $z$ 轴的角向平均（轴对称部分），$w=u-\mathcal{P}u$ 为其非轴对称部分。

\subsection{为什么“允许外力”不是文字游戏}

对不熟悉该问题的人，“允许外力”听起来像是把难题偷偷改简单了。要消除这种疑虑，得把费弗曼的措辞读细。

在备选 A 与 B 里，外力恒为零，命题针对\textbf{一切}容许初值作全称量化，要求解永远光滑；在备选 C 与 D 里，外力可以非零，但命题是存在性的——只要能指出\textbf{一组}容许的初值与外力，使得不存在一致有界能量的整体光滑解，就算完成。看似后者更弱，实则不然。困难恰恰来自对\textbf{外力本身}的约束：它必须在空间上速降（对 C，条件是 $\big|\partial_{x}^{\alpha}\partial_{t}^{m}f(x,t)\big|\le C_{\alpha mK}(1+|x|+t)^{-K}$ 对一切 $\alpha,m,K$ 成立；对 D，只需在时间上速降），也就是光滑、无奇性、可以取紧支。于是，把奇点“塞进外力”这条捷径被堵死了：解得靠自己发散，外力却必须始终是个规矩的函数。

这一条约束，把问题从“找出一条发散的解”改造成“找出一条发散的解，同时让它留下的残差规矩”。残差越大越好造，残差规矩越难造。2026 年工作的全部技术重量，几乎都压在后半句上。理解这一点，是理解后文所有标度、脉冲与修正的钥匙。

\begin{keybox}
\textbf{一句话}\quad 外力不是外挂的作弊器，它是问题的一部分，而且是被约束得最严的那一项。解题者被允许挑一个外力，却不被允许挑一个坏的外力。
\end{keybox}

\section{前史：反例系谱与真正的困难所在}

要理解 2026 年的构造，先要知道此前“反例系谱”已经长成什么样。这一节的用意不是罗列年表，而是标出困难的位置：一百年来，人们已经能造出各种各样的病态解，却一直造不出一个“外力光滑的”病态解。

\subsection{Leray 弱解与能量不等式}

1934 年，让·勒雷（Jean Leray）\cite{leray}证明三维 \NS{} 方程存在整体有限能量的弱解，并满足能量不等式
\begin{equation}
\|u(t)\|_{L^{2}}^{2}+2\nu\int_{0}^{t}\|\nabla u(s)\|_{L^{2}}^{2}\,ds\le\|u_{0}\|_{L^{2}}^{2}.
\end{equation}
弱解的存在性不成问题，成问题的是光滑性与唯一性。勒雷本人的构造允许解在有限时间内失去正则性，此后“再启动”为另一个弱解。因此，从勒雷起，真正未知的就是：光滑初值出发的那条解，能否一直保持光滑。

\subsection{CKN 部分正则性：奇点集不可能太大}

1982 年，Caffarelli、Kohn 与 Nirenberg\cite{ckn}证明“适当弱解”的奇点集在抛物尺度下的一维 Hausdorff 测度为零。这是关于奇点几何的第一条强约束：即便奇点存在，它们也只能坐落在一个极小的集合上。这条结果给后来的构造者出了一道隐含的题——如果要造奇点，最好把它造得足够“干净”，例如时间上孤立、空间上单点，否则会与 CKN 的约束冲突。

\subsection{条件正则性与临界范数}

Escauriaza、Seregin 与 Šverák\cite{ess}证明：若无外力 Cauchy 问题的解在 $L_{t}^{\infty}L_{x}^{3}$ 中一致有界，则解光滑。$L^{3}$ 恰好是方程在抛物伸缩下的临界空间，这类“临界范数有界即正则”的结果划出了爆破解必须越过的门槛：真出现爆破，速度必须在临界范数上无界。它同时告诉后来者，人为制造的爆破应当从“临界范数发散而低于临界的量保持有界”入手。

\subsection{Tao 的平均化方程：在简化模型里实现爆破}

Terence Tao\cite{tao-averaged}构造了一个经过平均化的 \NS{} 型方程，它保留原方程的非线性能量抵消，却允许在有限时间内爆破。这条结果的意义是方法论上的：它说明“爆破”本身不是不可企及的，真正卡住人的是原方程中那些被平均化抹去的强非线性项。此后人们开始系统地把爆破解的构造问题，转译为“如何在保留关键结构的前提下安排一种可控的不稳定性”。

\subsection{凸积分与非唯一性}

2019 年，Buckmaster 与 Vicol\cite{bv-nonunique}用凸积分方法证明有限能量弱解的非唯一性；随后 Albritton、Brué 与 Colombo\cite{abc}用相似变量下的不稳定涡旋，构造出初值为零、外力相同却彼此不同的 Leray--Hopf 解。后者尤其接近 2026 年的构造：它同样在一个“看起来应当光滑”的设定里制造出病态行为。差别在于，它的外力属于 $L_{t}^{1}L_{x}^{2}$，在初始时刻奇异；而千禧年问题的备选 C 要求外力在空间与时间上都光滑、速降。把外力的正则性从“奇异”提升到“光滑”，正是此后数年的主攻方向。

\begin{keybox}
\textbf{前史给出的关键判断}\quad 到 2020 年代中期，制造“某种意义下的爆破解”已经不难；难的是让构造过程中出现的那个外力不跟着一起变坏。换句话说，困难的坐标已经从“方程内部”平移到了“外力”这一项上。2026 年的工作，正是冲着这一项去的。
\end{keybox}

\subsection{Onsager 判据与能量级串：临界指数为何反复出现}

在前史里还有一条不能略过的线索，它是理解后文“尺度”与“可和性”为何处处卡在临界值的背景。Onsager 在 1949 年提出：三维湍流中，若速度场的 Hölder 指数 $s>1/3$，则能量守恒；若 $s<1/3$，则可能出现能量在有限时间内耗尽的“反常耗散”。$1/3$ 这个数正是使能量通量在尺度间守恒的临界指数。此后的数学工作把 Onsager 的猜想严格化：临界指数以上的解保持能量守恒（常数与 Eyink 等人），其下则存在能量非守恒的弱解（含本节的凸积分构造）。

这条线索与本文的关系有两层。第一层是术语上的：后文会反复出现“临界范数”“可和临界指数”“按尺度的通量”，这些都是同一个尺度的不同面孔，$L^{3}$、$B^{1/3}_{3,\infty}$、二维可和指数 $-3$，彼此只差一个由维度决定的常数。第二层是思路上的：爆破构造普遍依赖“能量在小尺度上不完全被黏性抹平”这一图景；而二次面（脉冲的 Reynolds 应力、涡层的剪切）正是从小尺度向大尺度回输能量的通道。2026 年 \NS{} 构造里“脉冲向背景补给动量”的安排，与湍流中“小尺度向大尺度回能”的逆级串，结构上是同一件事的两种说法。

\begin{principlebox}
\textbf{一个反复出现的数}\quad $1/3$、$L^{3}$、$-3$、$B^{1/3}_{3,\infty}$ 是同一个尺度在四套语言里的写法。凡是围绕该问题的技术文本，遇到这些数都应当本能地想到“临界”——它们标记的是“恰好能守恒、恰好能求和”的分界，而不是随手选的常数。
\end{principlebox}

\subsection{凸积分非唯一性与千禧年爆破是两件事}

前史里另有一处容易混淆的地方，值得单独澄清：\textbf{凸积分（convex integration）给出的非唯一性，与千禧年问题所问的爆破，是两件不同的事。}

凸积分证明的是“弱解不唯一”：同一初值（甚至同一外力）可以对应多个不同的有限能量弱解，其中一些能量不守恒。2019 年 Buckmaster--Vicol 的结果即属此类。它的要害在于解的正则性极低——弱到不满足能量等式，正是这种“低正则”给了构造自由度。

而千禧年备选 C 要求的是另一回事：要找一条在奇点前 $C^{\infty}$ 光滑、动能一致有界的解，让它在有限时间自发失去 $L^{\infty}$ 有界性。这里没有“弱解”的余地，也没有“能量等式失效”的余地。换言之，凸积分的自由来自“放弃正则性”，而 \NS{} 爆破的困难恰在于“不能放弃正则性”。

把这两个目标分开之后，就能看清 2026 年构造的定位：它既没有走弱解的路线，也没有走“数值找近似解再证明稳定”的路线，而是走了一条介于两者之间的第三条路——完全在光滑范畴内构造一个残差控制得住的奇异流。它的技术难度，正来自“既要完全光滑、又要真的发散”这一对要求本身。

\subsection{为什么二维已解决、三维未决：涡旋拉伸}

前史里还有一个最基本的问题没有正面回答：为什么同一组方程，二维的整体正则性早已被证明，三维却至今悬置？答案只有四个字——涡旋拉伸。

在二维，涡量 $\omega$ 是一个标量，沿流线守恒：$\partial_{t}\omega+u\cdot\nabla\omega=\nu\Delta\omega$。既然涡量既不能被拉伸、也不能被放大，临界范数自然保持有界，正则性随之而来（Clay 的官方提法把这一条归功于 Ladyzhenskaya 等人）。在三维，涡量是一个向量，方程里多出 $\omega\cdot\nabla u$ 这一项——涡管可以被流场拉长。拉长使涡量截面变细、强度变大，这一放大机制可能压过黏性的扩散效应，从而使临界范数失去控制。于是三维的问题不在于方程有何不同，而在于多了一个“可以把涡量越拉越强”的自由度。

这一观察为后文提供了两条线索。其一是机械性的：2026 年构造正是主动利用涡旋拉伸——核心沿轴向拉长、半径收缩，角向速度随之增大，与 $\omega\cdot\nabla u$ 的放大效应同气连枝。其二是概念性的：任何三维爆破解都必须让某个临界范数发散，而二维构造无论怎么设计都无法突破这条限制；因此 2026 年的证明在二维不可能有对应版本，这不是技术不足，而是维度本身设下的壁垒。

\section{思想脉络（一）：Córdoba--Martínez-Zoroa 的“残差博弈”}

\subsection{先选解，后定力：残差博弈}

设想一般的演化方程 $\partial_{t}U+\mathcal{N}(U)=f$。对任意光滑的 $U$，只要定义 $f:=\partial_{t}U+\mathcal{N}(U)$，$U$ 就自动是“带外力方程”的解。于是“存在一个使其爆破的外力”这句话在字面上是平凡的。全部内容都落在对 $f$ 的正则性要求上：$f$ 必须光滑、紧支，并连同一切导数穿过奇点。Diego Córdoba 与 Luis Martínez-Zoroa 在 2024 年的工作\cite{cmz-ipm}里把这件事说成了一场博弈：允许残差的各个分项发散，但不允许它们的和发散。后来的所有构造都在这条博弈规则下进行。

这解释了为什么这类证明往往冗长：不是因为想法复杂，而是因为要同时把“位移”“速度”“压力”“黏性”几项都调到恰好抵消的尺度上，其记账量极大。

\subsection{迭代放大与高阶修正}

Córdoba--Martínez-Zoroa 在二维不可压多孔介质方程（IPM）中给出了第一个完整样本\cite{cmz-ipm}：以光滑、紧支的源项，构造出有限能量的光滑解，在有限时间内产生奇点。其骨架包含两个相互咬合的机制。

其一是逐尺度放大。先在低频上造一个解，再加入一个高频、小幅的扰动；扰动在其线性化的演化方程中经历不稳定增长，把解推向更奇异的方向；然后在这个更奇异的解之上再叠一层更细的扰动。如此层层叠加，层与层之间的时间间隔可以求和为一个有限的 $T_{*}$，从而在有限时间里塞进无穷多层。这样，梯度与涡量在 $t\uparrow T_{*}$ 时发散。

其二是随层号提高修正阶数。每一层的高频 ansatz 在截断包络之外必然留下振荡误差；要保证外力光滑，就必须把误差抵消到足够高的阶。Córdoba 与 Martínez-Zoroa 的做法是让修正阶数 $J_{q}$ 随层号 $q$ 一起趋于无穷，使外力级数在每一个导数意义上都收敛。这条“阶数随尺度同步上升”的原则，是整套构造能给出 $C^{\infty}$ 外力而非仅有有限阶光滑的关键。

\subsection{从粗糙外力到光滑外力}

Córdoba 与 Martínez-Zoroa 的后续工作把 IPM 的思路推广到带耗散的 \NS{} 型方程与外力下的 Euler 方程\cite{cmz-euler}，但外力的正则性起初只达到局部适定类，还不是备选 C 所要求的光滑速降。把外力正则性推到 $C_{c}^{\infty}$，是留给下一棒的问题。

\subsection{Alpöge--Buckmaster 的推进}

2026 年 8 月，Levent Alpöge 与 Tristan Buckmaster 沿着这条纲领前进，在光滑外力下先后拿下 IPM、二维 Boussinesq 与三维不可压 Euler 的爆破，并以 Lean 完成验证。Buckmaster 在声明\cite{buckmaster-statement}中明确把纲领的归属记在 Córdoba 与 Martínez-Zoroa 名下，称该攻击路线的“基本想法”来自二人，甚至表示 Martínez-Zoroa 应当因此获得菲尔兹奖。Tao 在 9 月 7 日的博文\cite{tao-blog-sep7}中独立地给出了同样的溯源，并把这条路线概括为：在低频解上反复叠加高频小幅修正，使解在临近爆破时越来越奇异，同时保持外力良态。

对本文而言，重要的是这条路线的“角色分工”在 Alpöge--Buckmaster 那里出现了一次转变。在 Boussinesq 中，被放大的波本身构成奇点：每一层被指数放大的温度梯度，成为下一层更细波的仿射背景。而在 \NS{} 中，奇点是一个先写下来的光滑自相似坍缩，被放大的波退居为工具，用它们的平均动量通量去补足坍缩所缺的那份外力。这一分工的差异，正是第 6 章要展开的核心。

\subsection{Tao 的抽象迭代框架}

Tao 在同一篇博文里把上述思想抽象成一个不依赖具体方程的三段式。设已有低频（外力作用下的）解 $U_{0}$，欲在其上加一个高频小幅修正 $U=U_{0}+u$，代价是外力上多出一个修正 $g$。若能把 $u$ 的振幅做得相对大、而 $g$ 的振幅做得相对小，且让修正的频率随迭代次数迅速升高，就可以指望这一过程收敛到一个在有限时间爆破的解 $U$，而外力保持光滑。

Tao 进一步指出，要让“小扰动变大”，基本手段是让 $u$ 所满足的线性化方程
\begin{equation}
\partial_{t}u+\mathcal{L}u=g
\end{equation}
具有可利用的不稳定性：$u$ 可以在早期指数地小，而在临近爆破时被放大到足以破坏光滑性，同时又不至于大到使非线性分析失控。在 Boussinesq 里，这个线性化系统甚至有显式解——在原点附近速度关于空间线性、温度呈高频平面波——从而化为一组显式的振幅 ODE，其不稳定性可以直接读出。第 4 章将把这类“精确波”计算完整推一遍，因为它既是 Boussinesq 的引擎，也是 \NS{} 中脉冲生长与通量生成的理论依据。

\begin{principlebox}
\textbf{脉络小结}\quad Córdoba--Martínez-Zoroa 提供了纲领（残差博弈＋逐尺度放大＋阶数随尺度上升）；Alpöge--Buckmaster 把纲领推进到光滑外力与 Euler，并改变了被放大扰动的角色；Tao 把它抽象为“低频解＋可利用不稳定的高频修正”。OpenAI 的工作落在这一谱系的末端，其独有的贡献是把 \NS{} 的坍缩核心与“由离心不稳定性生长的脉冲所供给的 Reynolds 应力”精确地对接起来。
\end{principlebox}

\subsection{这条纲领的强项与短板}

把纲领摆到台面上，才能公正地评价 2026 年的推进。它的强项有三。其一，把“求解方程”转成了“设计抵消”，使问题从动力学变成可记账的代数，由此获得了极高的可操作性：一旦尺度对齐、阶数配好，剩下的就是逐项估计。其二，它对初值极不敏感，甚至可以从静止出发——因为解是被构造出来的，而初值只是一个待满足的边界条件。其三，它与外力正则性天然兼容：只要把修正阶数提到足够高，外力就能在任意多的导数意义下被压平。

它的短板同样清楚。第一，它给出的是存在性反例，而非对一般解的刻画：它说明“有些流会爆”，不说明“哪些流会爆”，更不触及无外力情形。第二，构造高度依赖人为设计，尺度、频率、修正阶数都需要精心选择，因而对“如果某一步的估计不够好会怎样”缺乏鲁棒性说明。第三，也是最容易被忽略的一点：这类构造的“可读性代价”很高——为了把外力压光滑，技术细节层层叠加，最终产物往往长达上百页，其核心思想被淹没在估计之中。第 7 章的逐阶修正与第 10 章的理解之争，追根溯源都来自这第三条短板。
